% Source: https://artofproblemsolving.com/community/c3h2655749_amc_10_problems
% Stable identity: aops-topic-2655749-q14
By Vieta's formulas the six roots (with multiplicity) have sum $10$ and product $16$. The product $16=2^4$ forces every root to be a power of $2$. Six powers of $2$ with product $2^4$ and sum $10$: the options are $\{1,1,2,2,2,2\}$ (sum $10$), $\{1,1,1,2,2,4\}$ (sum $11$), $\{1,1,1,1,4,4\}$ (sum $12$), $\{1,1,1,1,2,8\}$ (sum $14$) and $\{1,1,1,1,1,16\}$ (sum $21$). So the roots are $1,1,2,2,2,2$.

$B=-e_3$, where $e_3$ is the sum of all products of three roots:
\begin{itemize}
\item two $1$s and one $2$: $4$ products equal to $2$;
\item one $1$ and two $2$s: $2\cdot6=12$ products equal to $4$;
\item three $2$s: $4$ products equal to $8$.
\end{itemize}
So $e_3=8+48+32=88$ and $B=\boxed{-88}$.
